\documentclass[twocolumn]{article}
\usepackage[utf8]{inputenc}
\usepackage{amsmath}
\usepackage{hyperref}
\usepackage{graphicx}
\usepackage{caption}

\title{Optimización del problema ``Subir Escaleras'': de recursión exponencial a fórmula cerrada}
\author{Participantes: Jhon Perlaza, Heidi Sofía, Andrés Cadena, Juan Ahumada}
\date{}

\begin{document}
\maketitle

\begin{abstract}
Este trabajo estudia el problema clásico de contar cuántas formas existen de subir $n$ escalones usando pasos de tamaño 1 o 2. Se implementa una solución ingenua recursiva sin memoización, cuya complejidad temporal crece exponencialmente, y se compara contra una implementación optimizada basada en la relación con Fibonacci y la fórmula cerrada de Binet. Se presenta un análisis matemático de complejidad y un benchmarking que evidencia la diferencia de escalabilidad al aumentar $n$.
\end{abstract}

\section{Planteamiento del problema}
Sea $W(n)$ el número de formas de subir $n$ escalones usando pasos de $\{1,2\}$. El conteo cumple:
\[
W(n) = W(n-1) + W(n-2), \quad W(0)=1,\; W(1)=1.
\]
Esta recurrencia coincide con Fibonacci desplazado, ya que:
\[
W(n) = F(n+1).
\]

\section{Implementación pesada (exponencial)}
La versión pesada aplica recursión directa:
\[
W(n)=W(n-1)+W(n-2).
\]
Si $T(n)$ representa el tiempo de ejecución (número de llamadas), entonces:
\[
T(n)=T(n-1)+T(n-2)+O(1).
\]
Esta recurrencia tiene crecimiento exponencial, y se cumple:
\[
T(n)=\Theta(\varphi^n),
\]
donde $\varphi=\frac{1+\sqrt{5}}{2}$ es la proporción áurea. En particular, \textbf{no} es $O(n^2)$: el árbol de recursión se duplica aproximadamente con $n$ y el número de subproblemas crece como Fibonacci.

\section{Implementación optimizada (``O(1)'' por fórmula cerrada)}
Dado que $W(n)=F(n+1)$, puede usarse la fórmula cerrada de Binet:
\[
F(k)=\frac{\varphi^k-\psi^k}{\sqrt{5}},
\quad
\psi=\frac{1-\sqrt{5}}{2}.
\]
En el sentido de \emph{pasos matemáticos} (evaluar una fórmula sin iterar hasta $n$), esto se considera \emph{tiempo constante} $O(1)$. Sin embargo, en una implementación exacta con aritmética de alta precisión, el costo real depende del número de dígitos necesarios para representar $F(k)$ (modelo de bits). En este trabajo se usa un enfoque híbrido:
\begin{itemize}
\item Para $n$ pequeño: \texttt{float}, muy rápido y exacto en ese rango.
\item Para $n$ grande: \texttt{decimal} con precisión estimada para redondeo correcto.
\end{itemize}

\section{Benchmarking y resultados}
Se implementó un script de benchmarking que mide el tiempo para una lista fija de valores $n$. La versión exponencial se omite para $n>35$ para evitar tiempos excesivos. Se generan:
\begin{itemize}
\item \texttt{results/benchmark.csv}
\item \texttt{results/benchmark.png} (escala log--log)
\item \texttt{results/benchmark\_small.png} (zoom para $n\le 35$)
\end{itemize}

\begin{figure}[h]
\centering
\includegraphics[width=\linewidth]{results/benchmark_small.png}
\caption{Comparación en escala log--log para $n\le 35$.}
\end{figure}

\section{Aplicaciones reales}
La idea de ``reducir drásticamente'' complejidad mediante matemáticas, precomputación o técnicas ingeniosas aparece en:
\begin{itemize}
\item \textbf{Criptografía y teoría de números:} uso de exponenciación rápida ($O(\log n)$) y propiedades algebraicas para acelerar operaciones.
\item \textbf{Sistemas y bases de datos:} uso de índices y estructuras (B-trees, hash) para acceder en tiempo cercano a $O(1)$ u $O(\log n)$.
\item \textbf{Redes:} tablas de enrutamiento y estructuras compactas para decisiones rápidas.
\end{itemize}

\section{Repositorio}
Enlace del repositorio (pegar luego):
\[
\url{PEGAR_ENLACE_AQUI}
\]

\section{Referencias}
\begin{itemize}
\item Weisstein, E. W. ``Fibonacci Number.'' \emph{MathWorld}. \url{https://mathworld.wolfram.com/FibonacciNumber.html}
\item Wikipedia. ``Fibonacci number'' y ``Binet's formula''. \url{https://en.wikipedia.org/wiki/Fibonacci_number}
\item Knuth, D. E. \emph{The Art of Computer Programming}, Vol. 1 (discusión de recurrencias y crecimiento).
\end{itemize}

\end{document}